% !TeX root = ../papers/nc-p3.tex
\section{片側対称な階数十での全ての非局所自由境界}
\label{sec:finite-rank-ten-instanton-boundary}
\begin{theorem}\label{thm:finite-rank-ten-instanton-boundary}
有限点グラフを持つ片側対称な標準有限Frobenius窓で反対階数が10なら、
定理\ref{thm:finite-instanton-sheaf-boundary}の層が局所自由でない全ての窓は
二元三次式・二次式の乗法への辺商を持ち、非正則である。
従ってこの枝の全正則窓の添字保持同型類は、電荷4のインスタントン\emph{束}の
射影同値類に単射する。反対階数7から9の非局所自由窓には既存の全次数判定が適用される。
\end{theorem}
\begin{proof}[特異商上の二つの可換な一次鉛筆]
前定理の非零商 $Q$ を取り、一般の二つの一次式 $x,y$ による台の有限射影を
$f:\operatorname{Supp}_{\rm sch}Q\to\mathbb P^1$ とする。
$G=Q(-1)$ は純なので $f_*G$ は階数 $d$ のベクトル束であり、
$H^0(G(-1))=H^1(G(-1))=0$ により $f_*G=\mathcal O^d$ である。
残り二つの座標はこの自明束からその一次捩りへの可換な二行列
$A=xA_0+yA_1,B=xB_0+yB_1$ で作用する。
$H^0(G)\hookrightarrow H^1(\mathcal E(-1))=N_0$ は $H^0(F(-1))=0$ により単射で、
この埋込みは $H^0(G(1))\to N_1$ と乗法で両立する。
従って四行列の共通不変部分空間 $W\subset k^d$ への制限で成り立つ一次関係は、
全て $K\cap(U^*\otimes W)$ に入る。

次元が最小の非零共通不変部分空間を $W$、その次元を $n$ とする。
$n=1$ なら一つの $v$ に対し四座標の積が二次元以下に入るので、
二独立な一次式が $v$ を消し有限点グラフに反する。
$n\ge2$ の $W$ を同一の非零一次式が消すことも不可能である。いずれも $K$ に純関係でなく
点グラフの直線を与えるからである。
$n=2$ なら二行列の跡零部分は関数体上で比例する（非スカラー二次行列の中心化環は二次元）。
一方が零なら $W$ は一つの固定一次式で消える。ともに非零なら、成分の共通因子を除いた
原始行列が一次のときは比例係数が定数なので同じ矛盾になる。
原始行列が定数ならその固有直線が全係数の不変直線である。従って $n=3,4$ だけが残る。

$W$ が定める部分層のスキーム論的な台を $Y$ とする。
$f_*\mathcal O_Y$ は $\mathcal End(W\otimes\mathcal O)$ の部分代数なので、
直線束分解の次数は全て非正である。また大域切断は $A_0,A_1,B_0,B_1$ と可換な
定数行列であり、$W$ の既約性によりスカラーに限る。その階数 $e$ は $n$ 以下である。
実際二つの可換 $n$ 次行列の生成代数は次元 $n$ 以下である。
関数体の代数閉包に移り、$A$ のJordanブロックごとに異なる固有値を選んで、
$A$ と可換で巡回的な行列 $C$ を作る。一般の $B+tC$ は巡回的で $A$ と可換なので、
二行列の語はその $n$ 次元の中心化環に入る。$t=0$ で語の張る次元は増えない。
よって $f_*\mathcal O_Y=\mathcal O\oplus\bigoplus_{j=2}^e\mathcal O(-a_j)$、$a_j\ge1$ である。
\end{proof}

\begin{proof}[三次元部分と四次元部分]
$n=3$ なら四つの独立な一次座標の存在より $4\le h^0(\mathcal O_Y(1))\le2+(e-1)\le4$。
従って $e=3$、二つの $a_j$ は1である。
$Y$ は滑らかなねじれ三次曲線であることを示す。まず二つの非零一次式の積が
$Y$ を消すことはない。そうなら関数体上の二行列の積が零で、どちらも非零なので
その一つは階数1である。一次行列が関数体上で階数1なら、原始列・行に分けると
次数の和が高々1なので、像直線か核平面が定数である。
他の座標行列と可換だからこの定数部分空間は全係数で不変となり矛盾する。

$x,y$ と他の二座標は $f_*\mathcal O_Y(1)$ の全切断を張るので、
斉次座標環は $k[x,y]$ 上で次数 $0,1,1$ の自由基底を持つ。
実際他の二座標が二つの次数1の基底を与え、以後も二変数を掛けて全切断を張る。
そのイデアルは三二次式で生成される（他の二座標の三つの積を自由基底に戻す）。
高さ2のCohen--Macaulay環なのでHilbert--Burch分解とHilbert分子
$1-3z^2+2z^3$ により、この三式は二行三列の一次行列の小行列式である。
いずれかの非零の行・列結合の交点成分が零なら、その成分を左上に持ってきて非零小行列式に
二一次式の積を得る。三小行列式は独立なのでこれは先の禁止に反する。
よって全ての行結合は三独立な一次式を与え、その零点は一点である。
行方向の射影直線との入射は従って射影直線そのもので、次数3の小行列座標で写る。
三形式の共通の点は全て一意な行方向を持つ。行列全体がある点で零なら
全ての行方向の唯一の零点がそこに集中して曲線を覆えないからである。
単元成分を使って方向を局所的に解けるので入射から $Y$ への射はスキーム同型である。

$W$ に対応する $G$ の部分層はこの滑らかな三次上の階数1の捩れのない層で、
$\chi=3$ なので $\mathcal O_{\mathbb P^1}(2)$ である。
その三つの切断と四座標との積は $P_2\otimes P_3\to P_5$ の通常の乗法に等しい。
よって元の辺をこの三文字へ圧縮すると $R$ の像はこの乗法の双対像に入り、
定理\ref{thm:binary-linear-complex}の非正則な二元商となる。
ここで $G(1)$ から $N_1$ への写像が単射である必要はない。

$n=4$ なら $d=4$ かつ $W=k^4$ で、$N_0=H^0(G)$ である。
したがって二回乗法の作用は $H^0(\mathcal O_Y(2))$ を経由し、
その次元は高々 $3+2(e-1)\le9$ である。これは $\operatorname{Sym}^2 U^*\to
\operatorname{Hom}(N_0,N_2)$ の階数、すなわち反対階数が10であることに反する。
これで全ての非局所自由の場合が尽きる。
\end{proof}
